\section*{B Missing proofs from Section3}
\subsection*{B.1 Proofs from Section3.1}
\begin{authorclaim}{Claim8}
Let $P=(Q'_1)^{e_1}\cdots(Q'_t)^{e_t}$ be one of the polynomials $P_i$. For $k\ge0$, let $P^{(k)}:=\prod_{i\in[t]}(Q'_i)^{\max(e_i-k,0)}$. Then, $\partial_{\mathbf x}^kP\subseteq\mathbb F\text{-span}\{\mathbf y_M^{\le\infty}\mathbf x_M^{\le k(2t\tau-1)}P^{(k)}\}$.
\end{authorclaim}
\begin{proof}
We prove the claim by induction on $k$. If $k=0$, then $\partial_{\mathbf x_M}^0P=\{P\}=\{P^{(0)}\}$ and hence the claim is true. Assume that the claim is true for $k$. Let $X$ be a multilinear monomial of degree $k+1$ in $\mathbf x$ variables. Then $X=xX'$ where $X'$ is a multilinear monomial of degree $k$ in $\mathbf x$ variables and $x$ one of the $\mathbf x$ variables. From the induction hypothesis we have that,
\[
\frac{\partial P}{\partial X'}=g\cdot P^{(k)}
\]
where $g$ is a polynomial in $\mathbb F[\mathbf x_M,\mathbf y_M]$ with $\mathbf x_M$ degree of $g$ being at most $k(2t\tau-1)$ while its $\mathbf y_M$ degree can be arbitrarily large.
Let $J:=\{j\in[t]:e_j>k\}$. We have that,
\begin{align*}
\frac{\partial P}{\partial X}
&=\frac{\partial}{\partial x}(g\cdot P^{(k)})
=\frac{\partial}{\partial x}\left(g\cdot\prod_{j\in J}(Q'_j)^{e_j-k}\right)\\
&=\frac{\partial g}{\partial x}\prod_{j\in J}(Q'_j)^{e_j-k}
+g\sum_{j\in J}(e_j-k)(Q'_j)^{e_j-k-1}\frac{\partial Q'_j}{\partial x}
\prod_{i\in J\setminus\{j\}}(Q'_i)^{e_i-k}\\
&=\left(\frac{\partial g}{\partial x}\prod_{j\in J}Q'_j
+g\sum_{j\in J}(e_j-k)\frac{\partial Q'_j}{\partial x}
\prod_{i\in J\setminus\{j\}}Q'_i\right)
\prod_{j\in J}(Q'_j)^{e_j-k-1}.
\end{align*}
Observe that as $\mathtt D$ is a pruned depth four circuit, the support of all monomials of $Q'_j$ is upper bounded by $\tau$ and as in any monomial the individual degree of any $\mathbf x$ variable is at most two, $\deg_{\mathbf x}(Q'_j)\le2\tau$. Also, $|J|\le t$ and hence
\[
\deg_{\mathbf x}\left(\frac{\partial g}{\partial x}\prod_{j\in J}Q'_j
+g\sum_{j\in J}(e_j-k)\frac{\partial Q'_j}{\partial x}
\prod_{i\in J\setminus\{j\}}Q'_i\right)\le(k+1)(2t\tau-1).
\]
As $\prod_{j\in J}(Q'_j)^{e_j-k-1}=P^{(k+1)}$, the claim is true for $k+1$.
\end{proof}
\begin{authorclaim}{Claim9}
Let $\ell,k,t$ and $\tau$ be as defined earlier. Then, $\ell+2kt\tau<m/2$.
\end{authorclaim}
\begin{proof}
We will show that the ratio $(m/2-2kt\tau)/\ell>1$. Putting the values of $k$ and $\ell$,
\[
\frac{m/2-2kt\tau}{\ell}
=\frac{m/2-2\lfloor\delta d_x/t\rfloor t\tau}{\lfloor m/(m^{\delta/t}+1)\rfloor}
\ge\left(\frac12-\frac{2\delta d_x\tau}{m}\right)(m^{\delta/t}+1).
\]
So, we need to show that
\begin{align*}
\frac1{1/2-2\delta d_x\tau/m}<m^{\delta/t}+1
&\Longleftrightarrow\frac1{1/2-2\delta d_x\tau/m}-1<m^{\delta/t}\\
&\Longleftrightarrow\frac{1+4\delta d_x\tau/m}{1-4\delta d_x\tau/m}<m^{\delta/t}.
\end{align*}
For large enough $m$, $4\delta d_x\tau/m\le1/2$. Using $1+x\le e^x$, which holds for all $x\in\mathbb R$, and $1/(1-x)\le e^{2x}$, which holds for $0\le x\le1/2$ we get:
\[
\frac{1+4\delta d_x\tau/m}{1-4\delta d_x\tau/m}\le e^{12\delta d_x\tau/m}.
\]
So showing that $e^{12\delta d_x\tau/m}<m^{\delta/t}$ would suffice. Now,
$e^{12\delta d_x\tau/m}<m^{\delta/t}\Longleftrightarrow e^{12d_xt\tau/m}<m$.
Putting the values of $d_x,t$ an $\tau$, we get that
\[
\frac{12d_xt\tau}{m}
=\frac{12d_x\lfloor d_x/(\ln m)^3\rfloor\lfloor20\ln m\rfloor}{m}
\le\frac{12d_x^2\cdot20\ln m}{m(\ln m)^3}
=\Theta\left(\frac m{m(\ln m)^2(\ln m)^2}\right)
=\Theta\left(\frac1{(\ln m)^4}\right)=o(1)
\]
as $d_x=\Theta(\sqrt m/\ln m)$. Thus $e^{12d_xt\tau/m}<m$.
\end{proof}
\subsection*{B.2 Proof from Section3.2.1}
\begin{authorclaim}{Claim11}
Procedure1 terminates in at most $m$ iterations.
\end{authorclaim}
\begin{proof}
Let $H_i$ be the set $H$ after the $i$-th iteration of the procedure. Since each monomial in $H_i$ has support more than $\tau$, for any such monomial there are at least $\tau/2$ distinct $j\in[3m]\setminus M_1$ such that at least one of $x_j$ and $y_j$ appears in it. Counting the number of times at least one of $x_j$ and $y_j$ appears in a monomial in $H_i$ and summing up these counts for all $j\in[3m]\setminus M_1$, we get that
\[
\sum_{j\in[3m]\setminus M_1}e(j)\ge\tau|H_i|/2;
\]
so from an averaging argument there exists a $j$ such that $e(j)\ge\tau|H_i|/(6m)$.
Hence, the size of $H_{i+1}$ is upper bounded as $|H_{i+1}|\le|H_i|(1-\tau/(6m))$.
So after $i$ iterations of the procedure we get,
\begin{align*}
|H_i|&\le|H_0|(1-\tau/(6m))^i\\
&\le\lfloor m^2d_x/(\ln m)^5\rfloor(1-\lfloor20\ln m\rfloor/(6m))^i\\
&\le\frac{m^2d_x}{(\ln m)^5}(1-(20\ln m-1)/(6m))^i\\
&\le\frac{m^2d_x}{(\ln m)^5}e^{-3i\ln m/m}\quad(\text{for sufficiently large }m)\\
&=\frac{m^2d_x}{(\ln m)^5}m^{-3i/m}.
\end{align*}
For $i=m$, $|H_i|<1$ (for sufficiently large $m$), i.e., the procedure terminates in at most $m$ iterations.
\end{proof}
\subsection*{B.3 Proof from Section3.2.2}
\begin{authorclaim}{Claim13}
Let $\overline M_1=[3m]\setminus M_1$. Procedure2 sets at most $m$ many variables in $\mathbf x_{\overline M_1}\cup\mathbf y_{\overline M_1}$ to field constants and removes all the heavy gates from $\mathtt C_1$.
\end{authorclaim}
\begin{proof}
In each iteration, we evaluate a light sparse polynomial in $\mathtt C_1$ to zero. This can be done as $\mathbb F$ is an algebraically closed field. Since the support of every monomial in $\mathtt C_1$ is at most $\tau$, we end up setting at most $\tau m/(\ln m)^2\le20m/\ln m$ many variables to field constants in each iteration. As we can afford to set $m$ variables, Step2 of the procedure executes successfully. For some $i\in\mathbb N$, the while loop terminates in the $i$-th iteration in either of the following two cases:
\begin{enumerate}
\item All the heavy gates get eliminated after the $(i-1)$-th iteration, i.e., $s_i=0$.
\item $\tau(b_1+\cdots+b_i)>m$. (We show in the following subclaim that all the heavy gates are eliminated before this happens. Hence, the procedure stops only in the above case.)
\end{enumerate}
\begin{authorclaim}{Subclaim18}
Let $i\in\mathbb N$ be such that $\tau(b_1+\cdots+b_{i-1})\le m$ but $\tau(b_1+\cdots+b_i)>m$. Then, all the heavy gates in $\mathtt C_1$ get eliminated in the first $(i-1)$ iterations of Procedure2.
\end{authorclaim}
If we assume Subclaim18 then Claim13 is proved.
\end{proof}
\begin{proof}[Proof of Subclaim18]
For $1\le j<i$, let $(Q_{j,1},\ldots,Q_{j,r_j})$ be the available light sparse polynomials in $\mathtt C_1$ after the $(j-1)$-th iteration. Recall that $s_j$ is the number of heavy gates in $\mathtt C_1$ after the $(j-1)$-th iteration. Suppose, $s_j\ge1$ (otherwise, we have nothing to prove). For every $l\in[r_j]$, $b_{j,l}$ and $c_{j,l}$ refer to the fan-in of $Q_{j,l}$ and the number of distinct heavy gates connected to $Q_{j,l}$ in $\mathtt C_1$ respectively. We may assume that $b_{j,l}\ne0$ for every $l\in[r_j]$. It is given that
\begin{equation}
b_{j,1}+\cdots+b_{j,r_j}\le\frac{m^2d_x}{160\lambda_0(\ln m)^5}.\tag{2}
\end{equation}
Since every heavy gate is connected to at least $md_x/(2\lambda_0(\ln m)^3)$ many light sparse polynomials in $\mathtt C_1$,
\[
s_j\frac{md_x}{2\lambda_0(\ln m)^3}\le c_{j,1}+\cdots+c_{j,r_j}.
\]
As $b_{j,1},\ldots,b_{j,r_j}$ are all non-zero, we get
\[
s_j\frac{md_x}{2\lambda_0(\ln m)^3}
\le\frac{c_{j,1}}{b_{j,1}}b_{j,1}+\cdots+\frac{c_{j,r_j}}{b_{j,r_j}}b_{j,r_j}.
\]
Let $u\in[r_j]$ be such that $c_{j,u}/b_{j,u}=\max\{c_{j,1}/b_{j,1},\ldots,c_{j,r_j}/b_{j,r_j}\}$. Let $c_j:=c_{j,u}$ and $b_j:=b_{j,u}$. Then, the above equation implies
\[
s_j\frac{md_x}{2\lambda_0(\ln m)^3}\le\frac{c_j}{b_j}(b_{j,1}+\cdots+b_{j,r_j}).
\]
From Equation(2), we get that for every $1\le j<i$,
\[
s_j\frac{md_x}{2\lambda_0(\ln m)^3}
\le\frac{c_j}{b_j}\frac{m^2d_x}{160\lambda_0(\ln m)^5},
\]
which implies
\begin{equation}
\frac{80s_jb_j(\ln m)^2}{m}\le c_j.\tag{3}
\end{equation}
Thus, by setting the light sparse polynomial $Q_{j,u}$ to zero in the $j$-th iteration, we get rid of at least $80s_jb_j(\ln m)^2/m$ many heavy gates from $\mathtt C_1$. Recall that $s_{j+1}$ is the number of available heavy gates after the $j$-th iteration. Then, for every $1\le j<i$,
\begin{equation}
s_{j+1}\le s_j-c_j\le s_j(1-80b_j(\ln m)^2/m).\tag{4}
\end{equation}
Hence, for every $1\le j<i$,
\begin{equation}
s_{j+1}\le s_1\prod_{l=1}^j(1-80b_l(\ln m)^2/m).\tag{5}
\end{equation}
Also, for every $l\le j$, we have $c_l\le s_l$. Thus, Equation(3) implies
\begin{equation}
80b_l(\ln m)^2/m\le1.\tag{6}
\end{equation}
As $1-a\le e^{-a/2}$ for $0\le a\le1$, Equations(5) and(6) imply
\begin{equation}
s_{j+1}\le s_1\prod_{l=1}^je^{-40b_l(\ln m)^2/m}
=s_1e^{-40(b_1+\cdots+b_j)(\ln m)^2/m}.\tag{7}
\end{equation}
It is given that $\tau(b_1+\cdots+b_i)>m$, which implies $\tau(b_1+\cdots+b_{i-1})>m-\tau b_i$. As $b_i$ is the fan-in of a light sparse polynomial, $b_i\le m/(\ln m)^2$ and so
\begin{equation}
\tau(b_1+\cdots+b_{i-1})>m-\frac{m\tau}{(\ln m)^2}.\tag{8}
\end{equation}
On substituting $j=i-1$, $\tau=\lfloor20\ln m\rfloor$ and the value of $\tau(b_1+\cdots+b_{i-1})$ from Equation(8) in Equation(7), we get
\[
s_i\le s_1\exp\left(-\frac{40(\ln m)^2}{m}
\left(\frac m{20\ln m}-\frac m{(\ln m)^2}\right)\right)
=s_1e^{-(2\ln m-40)}.
\]
For large enough $m$, $s_i\le s_1/m^{1.9}$. Since $s_1\le m$, we get $s_i=0$ (as it is a natural number). In other words, we get rid of all the heavy gates within $(i-1)$ iterations.
\end{proof}
